% Options for packages loaded elsewhere
\PassOptionsToPackage{unicode}{hyperref}
\PassOptionsToPackage{hyphens}{url}
\PassOptionsToPackage{dvipsnames,svgnames,x11names}{xcolor}
\documentclass[
  10pt,
  fontset=fandol]{ctexart}
\usepackage{xcolor}
\usepackage[margin=16mm]{geometry}
\usepackage{amsmath,amssymb}
\setcounter{secnumdepth}{-\maxdimen} % remove section numbering
\usepackage{iftex}
\ifPDFTeX
  \usepackage[T1]{fontenc}
  \usepackage[utf8]{inputenc}
  \usepackage{textcomp} % provide euro and other symbols
\else % if luatex or xetex
  \usepackage{unicode-math} % this also loads fontspec
  \defaultfontfeatures{Scale=MatchLowercase}
  \defaultfontfeatures[\rmfamily]{Ligatures=TeX,Scale=1}
\fi
\usepackage{lmodern}
\ifPDFTeX\else
  % xetex/luatex font selection
\fi
% Use upquote if available, for straight quotes in verbatim environments
\IfFileExists{upquote.sty}{\usepackage{upquote}}{}
\IfFileExists{microtype.sty}{% use microtype if available
  \usepackage[]{microtype}
  \UseMicrotypeSet[protrusion]{basicmath} % disable protrusion for tt fonts
}{}
\makeatletter
\@ifundefined{KOMAClassName}{% if non-KOMA class
  \IfFileExists{parskip.sty}{%
    \usepackage{parskip}
  }{% else
    \setlength{\parindent}{0pt}
    \setlength{\parskip}{6pt plus 2pt minus 1pt}}
}{% if KOMA class
  \KOMAoptions{parskip=half}}
\makeatother
\setlength{\emergencystretch}{3em} % prevent overfull lines
\providecommand{\tightlist}{%
  \setlength{\itemsep}{0pt}\setlength{\parskip}{0pt}}
\usepackage{amsmath,amssymb,mathtools,needspace}
\xeCJKDeclareCharClass{CJK}{"2032,"0400->"04FF,"2160->"216B,"2460->"2473,"25A0,"25C7,"25CB,"25CF}
\IfFontExistsTF{Arial Unicode MS}{\xeCJKsetup{AutoFallBack=true}\setCJKfallbackfamilyfont{\CJKrmdefault}{Arial Unicode MS}}{}
\setcounter{secnumdepth}{0}
\setlength{\emergencystretch}{2em}
\allowdisplaybreaks
\usepackage{bookmark}
\IfFileExists{xurl.sty}{\usepackage{xurl}}{} % add URL line breaks if available
\urlstyle{same}
\hypersetup{
  pdftitle={相对误差与相对误差限},
  colorlinks=true,
  linkcolor={Maroon},
  filecolor={Maroon},
  citecolor={Blue},
  urlcolor={Blue},
  pdfcreator={LaTeX via pandoc}}

\title{相对误差与相对误差限}
\author{}
\date{}

\begin{document}
\maketitle

\paragraph{1.3.2 相对误差与相对误差限}\label{ux76f8ux5bf9ux8befux5deeux4e0eux76f8ux5bf9ux8befux5deeux9650}

误差限的大小还不能完全表示近似值的好坏. 例如, 有两个量 \(x = 10 \pm 1, y = 1000 \pm 5\), 则

\[
x^* = 10, \varepsilon_x^* = 1, y^* = 1000, \varepsilon_y^* = 5.
\]

虽然 \(\varepsilon_y^*\) 比 \(\varepsilon_x^*\) 大 4 倍, 但 \(\frac{\varepsilon_y^*}{y^*} = \frac{5}{1000} = 0.5\%\) 比 \(\frac{\varepsilon_x^*}{x^*} = \frac{1}{10} = 10\%\) 要小得多, 这说明 \(y^*\) 近似 \(y\) 的程度比 \(x^*\) 近似 \(x\) 的程度要好得多. 所以, 除考虑误差的大小外, 还应考虑准确值 \(x\) 本身的大小. 近似值的误差 \(e^*\) 与准确值 \(x\) 的比值

\[
\frac{e^*}{x} = \frac{x^* - x}{x}
\]

称为近似值 \(x^*\) 的相对误差, 记作 \(e_r^*\).

在实际计算中, 由于真值 \(x\) 总是不知道的, 通常取

\[
e_r^* = \frac{e^*}{x^*} = \frac{x^* - x}{x^*}
\]

作为 \(x^*\) 的相对误差, 条件是 \(e_r^* = \frac{e^*}{x^*}\) 较小, 此时

\[
\frac{e^*}{x} - \frac{e^*}{x^*} = \frac{e^* (x^* - x)}{x^* x} = \frac{(e^*)^2}{x^* (x^* - e^*)} = \frac{(e^* / x^*)^2}{1 - (e^* / x^*)}
\]

是 \(e_r^*\) 的二次方项级, 故可忽略不计.

相对误差也可正可负, 它的绝对值上界叫做相对误差限, 记作 \(\varepsilon_r^*, \varepsilon_r^* = \frac{\varepsilon^*}{|x^*|}\).

根据定义, \(\frac{\varepsilon_x^*}{|x^*|} = 10\%\) 与 \(\frac{\varepsilon_y^*}{|y^*|} = 0.5\%\) 分别为 \(x\) 与 \(y\) 的相对误差限, 可见 \(y^*\) 近似 \(y\) 的程度比 \(x^*\) 近似 \(x\) 的程度好.

\end{document}
